%  FRAM: the same catalogue idea as the 24C EEPROM, on I2C and SPI, without the
%  page boundary or the program cycle.
\chapter[FRAM: Non-Volatile RAM]%
        {FRAM: Non-Volatile RAM --- the Same Catalogue, Without the Wait}
\label{ch:fram}

FRAM (ferroelectric RAM) fills the same kilobytes-of-non-volatile niche as the 24C
EEPROM (Chapter~\ref{ch:eeprom}), and its driver is built the same way --- a
generic over a catalogue of parts. So this chapter is mostly about what FRAM does
\emph{not} need.\index{FRAM}\index{ferroelectric RAM}

It behaves like RAM that keeps its contents: byte-writable, effectively unlimited
endurance ($10^{12}$ writes and more), and no erase. Two of the three EEPROM traps
simply vanish. There is no page boundary, so a \texttt{Write} of any length is one
transaction with no split; and there is no program cycle --- a byte is committed
as it is clocked in --- so there is no 5\,ms wait and no ACK polling. FRAM's
\texttt{Write} is the EEPROM's with the two hardest parts deleted. What survives is
the addressing.

\section{Two buses, two families}

FRAM comes on I2C (Fujitsu MB85RC, Cypress FM24) and on SPI (Fujitsu MB85RS,
Cypress FM25), so there are two catalogues, \texttt{ESP32S3.FRAM\_I2C} and
\texttt{ESP32S3.FRAM\_SPI}, each a generic over its own \texttt{Geometry}. The I2C
side reuses the 24C addressing wholesale --- device-type code 1010, a big-endian
word address, and high address bits folded into the device-select byte on the
small parts, with the same derived \texttt{Blocks} / \texttt{Max\_Devices} and the
same per-instance strap precondition (Chapter~\ref{ch:eeprom}). The SPI side is the
small standard SPI-memory command set: \texttt{WREN} (0x06) to arm a write, then
\texttt{WRITE} (0x02) / \texttt{READ} (0x03) with a big-endian address, and
\texttt{RDID} (0x9F) for a JEDEC-style identity.

\section{Keyed by density, not by part number}

The read/write protocol is identical across manufacturers --- a fact the
datasheets made emphatic --- so a part is nothing but its geometry, and the
catalogue names instances by density rather than by vendor SKU. One
\texttt{Kbit\_256} covers the Fujitsu MB85RC256V and the Cypress FM24W256 alike;
its banner lists the part numbers.

\begin{lstlisting}
--  I2C FRAM
with ESP32S3.FRAM_I2C.Kbit_256;
...
   package Ram_Part renames ESP32S3.FRAM_I2C.Kbit_256;   --  MB85RC256V, FM24W256, ...
   Ram : Ram_Part.Device;
begin
   Ram_Part.Setup (Ram, Sda => 41, Scl => 40);   --  clock defaults per geometry
   Ram_Part.Write (Ram, 16#0100#, Buf, St);       --  one transaction, no split
\end{lstlisting}

The I2C catalogue runs \texttt{Kbit\_4} through \texttt{Mbit\_1} (512 bytes to
128\,KiB); the SPI catalogue \texttt{Kbit\_4} through \texttt{Mbit\_4}. Each
geometry also carries its own default bus clock --- FRAM runs faster than a 24C
part, up to Fast-mode Plus on I2C --- which \texttt{Setup} uses as the default for
\texttt{Clock\_Hz}.

The SPI parts take their pins and an application-driven chip select, exactly like
the W25Q flash (Chapter~\ref{ch:storage}): a plain \texttt{CS\_Pin}, or
\texttt{No\_Pin} plus a \texttt{CS\_CB} callback for a select behind an I/O
expander. Because SPI has no per-transaction ACK and FRAM has no program cycle to
fail, the SPI \texttt{Read} / \texttt{Write} carry \emph{no} \texttt{Status} at
all --- an absent chip shows up as a bad Device ID, not as a failed transfer.

\begin{lstlisting}
--  SPI FRAM
with ESP32S3.FRAM_SPI.Kbit_256;
...
   package Ram_Part renames ESP32S3.FRAM_SPI.Kbit_256;
   Ram : Ram_Part.Device;
begin
   Ram_Part.Setup (Ram, Sclk => 1, Mosi => 4, Miso => 45, CS_Pin => 12);
   Ram_Part.Write (Ram, 16#0100#, Buf);           --  WREN, then WRITE + payload
\end{lstlisting}

\section{Identify, do not verify}

Because FRAM can report a Device ID, it is tempting to read it back and confirm the
populated part matches the geometry you compiled for. The datasheets killed that
idea. On I2C the ID is a 3-byte reply --- a 12-bit manufacturer code (Fujitsu
\texttt{0x00A}, Cypress \texttt{0x004}), a vendor-specific density code, and an
8-bit product code --- read from the reserved slave address \texttt{0xF8}; on SPI
it is the JEDEC-style \texttt{RDID} reply. Either way the command is \emph{not}
implemented on most parts --- the manufacturer datasheets confirm the 16 / 64 /
128\,Kbit MB85RC parts carry none --- so a part that reports nothing is the common
case, not an error.

The driver therefore keeps \texttt{Read\_Device\_ID} and \texttt{Identify}
(\texttt{Fujitsu} / \texttt{Cypress} / \texttt{Unknown}) as \emph{informational}
only; the geometry is fixed at compile time by the chosen instance, and trusted. As
with the 24C parts, every FRAM geometry is datasheet-derived and marked
\texttt{STATUS: UNTESTED} until it is run on real silicon.

\noindent The on-target example \texttt{esp32s3\_fram} drives an I2C
MB85RC256V --- it reads the Device ID, writes a 100-byte pattern in one
transaction (past the I2C FIFO, to prove there is no split), and keeps the same
non-volatile boot counter the EEPROM demo uses. A companion \texttt{Fram\_Coverage}
unit \texttt{with}s every instance in both families, so the whole catalogue is
build-checked even though only one part is driven.

\vfill
\noindent\rule{\linewidth}{0.4pt}\\[4pt]
\textit{The whole stack (runtime, peripheral HAL, and filesystem) is one coherent
body of Ada: strong types from the register up, the same RAII and task-safety
patterns in every driver, and a filesystem validated against the Linux kernel, all
running on both cores of the ESP32-S3 on its own bare-metal runtime.}
